\section{Introduction}
\label{sec:intro}

% =====================================================================
% Block 1 — Network reconstruction as a core problem in complex systems
% =====================================================================
Many natural phenomena are usefully described as complex systems: large collections of
components whose collective behaviour emerges from a web of nonlinear interactions
\citep{strogatz2001exploring,barzel2013universality}. A central goal in their study is
to recover the directed network of interactions that governs how the components of such
a system influence one another \citep{barzel2013universality}. The same task recurs
across domains --- gene regulatory networks in molecular biology, signalling cascades in
cell biology, species-interaction networks in community ecology, and reaction--diffusion
processes in chemical systems \citep{timme2014revealing,turing1952chemical} --- because
these systems share a common mathematical core: each node evolves according to a rate
law driven by the states of the nodes that act on it, so the network is the support of
those driving terms. Recovering that support from observations of the system's dynamics,
rather than measuring it directly, is fundamentally an inverse problem
\citep{timme2014revealing,gao2022autonomous}. In microbial ecology, the
structure of the interaction network among taxa shapes community assembly, stability,
and resistance to invasion, and is increasingly viewed as a target for rational
therapeutic design \citep{clark2021design,ishizawa2024beyond}. For example,
modelling disrupted gut communities has identified protective species, such as
\textit{Clostridium scindens}, whose reintroduction can restore colonization resistance
against \textit{C. difficile}, a major cause of antibiotic-associated diarrhoea
\citep{buffie2015precision}.

% =====================================================================
% Block 2 — ODE modelling of dynamical networks: the parametric dilemma
% =====================================================================
Ordinary differential equations (ODEs) provide the natural language for this core: each
component's rate of change is a function of the current state, so the interaction
network is encoded directly in the right-hand-side functions, and reconstructing it
reduces to learning those functions and their support from data --- a goal driving rapid
recent progress at the interface of dynamical-systems modelling and machine learning
\citep{karniadakis2021physics,gao2022autonomous}. Two modelling choices then determine
what can be recovered. The first is the functional form of each interaction. A rigid
parametric form --- in the limit a single coefficient per interaction --- is
interpretable but cannot represent effects that saturate, switch sign, or act
nonlinearly, even though qualitatively distinct system-level behaviour follows from the
\emph{form} of the local interaction functions, not merely their sign
\citep{barzel2013universality}. Nonparametric additive ODE models avoid this commitment:
each node's rate is a sum of univariate functions of the other nodes, replacing
coefficients with smooth components while preserving an interpretable edge-wise
decomposition \citep{henderson2014network,wu2014saode,chen2017network}. The second choice is
regularization. Because real interaction networks are sparse --- most components are
directly influenced by only a few others, often with heavy-tailed connectivity
\citep{barabasi1999emergence} --- imposing sparsity on which univariate components are
nonzero turns estimation into network recovery. Coupling nonlinear additive ODEs with
group- and element-level sparsity scales this idea to genome-wide networks
\citep{dong2026multitask}, making sparse additive ODEs a practical route to large,
nonlinear, directed network inference.

% =====================================================================
% Block 3 — Data paradigms: time-series vs perturbed steady state
% =====================================================================
The prevailing paradigm in such ODE modelling is time-series inference: trajectories
are sampled over time, the derivative of each state is estimated, and these estimated
rates are regressed on candidate driving terms. This underlies both generalized
Lotka--Volterra (gLV) engines such as MDSINE \citep{bucci2016mdsine} and library-based
discovery of nonlinear dynamics, as in
sparse identification of nonlinear dynamics (SINDy)
\citep{brunton2016sindy,kaptanoglu2022pysindy}, including recent applications to
microbial pairwise interactions \citep{sindy2023microbiota}. Such approaches are
powerful when high-quality trajectories are available, but this requirement is often
the limiting factor in biological and ecological systems. Accurate derivative
estimation requires dense sampling during transient regimes, whereas experiments
typically provide noisy, irregular, and sparse abundance measurements; in microbial
communities, repeated high-density sampling can be invasive, costly, or simply
infeasible, and the number of interacting taxa can far exceed the number of observed
time points \citep{stein2013ecological,bucci2016mdsine}. Pseudo-time approaches partly
relax the need for explicit time stamps by ordering snapshot measurements along an
inferred trajectory --- for instance a principal or diffusion coordinate --- so as to
emulate a time course \citep{trapnell2014pseudotime,haghverdi2016diffusion}, after which
dynamics and regulatory networks can be inferred with ODE-based methods that treat
pseudotime as time \citep{matsumoto2017scode}. Yet the inferred ordering is itself
uncertain and strongly method-dependent \citep{saelens2019comparison}, and the
downstream step still reduces to estimating rates along a reconstructed trajectory. The
resulting numerical-differentiation problem is ill posed and amplifies measurement
error, making time-series-based network recovery fragile precisely in the low-budget
settings where network inference is most needed.

These limitations motivate a complementary paradigm that sidesteps the derivative
entirely by working at equilibrium. In microbial ecology, steady-state abundance data
alone can reveal network topology and interaction signs without assuming a dynamics
model, recovering interaction strengths only once a fixed (e.g. gLV) form is imposed
\citep{xiao2017mapping}. More broadly, modular response analysis recovers the topology
and strength of network connections from the steady-state responses of a system to
successive perturbations of its components
\citep{kholodenko2002untangling,jimenezdominguez2021aimera}, and scales to systems with
hundreds of nodes \citep{mekedem2022mra}. However, these methods characterise
interactions only locally or under a fixed parametric form, yielding a linearised
interaction matrix rather than the underlying nonlinear couplings. The perturbed
steady-state (\pss{}) approach of \citet{meister2013pss} for
gene regulation makes the key structural observation that, by restricting attention to
perturbed steady states where $dx/dt = 0$, a nonlinear ODE becomes an \emph{algebraic}
system in the unknown interaction functions, removing the need to estimate time
derivatives while retaining the nonlinear model. To date this reformulation has been
exploited only with a fixed parametric (biophysics-based) rate law, leaving the
nonparametric, network-scale setting open.

% =====================================================================
% Block 4 — PSS-Net: nonparametric additive ODE with double sparsity
% =====================================================================
We propose \pss{}-Net to fill this gap. The central idea is to combine the algebraic
\pss{} reformulation with a sparse additive ODE model: each directed interaction is
represented by a univariate basis expansion, and the steady-state equations turn
network recovery into a structured sparse-regression problem rather than a
trajectory-differentiation problem. Sparsity is imposed at two levels, selecting both
which source nodes regulate a target and which basis terms are needed to describe each
retained edge function. We solve this double-sparse problem using the \adsiht{}
estimator \citep{zhang2023minimax,cai2022sparsegroup}. In this sense, \pss{}-Net is
closest in spirit to MRA because both exploit perturbation-induced steady-state
responses, but it aims beyond a local linear response matrix by recovering sparse
nonlinear edge functions and their local Jacobian signs. Thus, \pss{}-Net inherits the
interpretable edge-wise structure of additive ODE network models
\citep{henderson2014network,wu2014saode,chen2017network}, avoids numerical
differentiation, and
generalises parametric \pss{} inference \citep{meister2013pss} to nonlinear,
network-scale recovery.

% =====================================================================
% Block 5 — Contributions
% =====================================================================
\paragraph{Contributions.}
\begin{itemize}
  \item We show that perturbed steady-state observations can support accurate recovery
        of sparse nonlinear interaction networks without numerical differentiation,
        and characterise how recovery depends on sample size, model nonlinearity, and
        perturbation design.
  \item We benchmark \pss{}-Net against MRA-style, sparse-regression, SINDy-style, and
        association-based alternatives, identifying regimes where nonlinear
        steady-state modelling improves support recovery and where topology or
        measurement constraints remain limiting.
  \item We reserve the final application layer for real microbial-community data, where
        \pss{}-Net can be used to analyse perturbation-response experiments once
        suitable steady-state abundance measurements are available.
\end{itemize}
